\chapter{Estado del arte}
\label{cap:sota}

% =====================================================================================
% GUION. Reescrito por completo en sep-2026.
%
% ⚠️ EL BORRADOR ANTERIOR ESTABA LLENO DE MARCADORES «[cite: 1, 2]» que NO son citas de
%    LaTeX: venian de una exportacion y se imprimian tal cual en el PDF. Sustituidos por
%    \cite{} de verdad.
%
% 🔴 RESTRICCION DURA DE EXTENSION. Los capitulos 2 y 3 son el bloque «Contexto y estado
%    del arte» y la norma lo acota a 10-15 PAGINAS EN TOTAL. El capitulo 2 ocupa 12, asi
%    que este debe caber en TRES. De ahi las decisiones de economia:
%      - cada trabajo se despacha en dos o tres frases, no en un parrafo;
%      - los paradigmas alternativos (Q-Cluster, Liao 2025) van en UNA frase dentro de
%        §3.2, no en seccion propia; su sitio es el capitulo de limitaciones;
%      - §3.5 va con TABLA: ahorra dos paginas de prosa y es justo lo que un tribunal
%        busca en un trabajo de tipo 3.
%    Si el capitulo 2 adelgaza, lo primero que merece crecer aqui es §3.3.
%
% Fuentes: TFM-Quantum/doc/SoTA/comparative_analysis.md (corpus principal) y
%          TFM-Quantum/doc/SoTA_ampliada/fichas/paper_{1..5}.md (solo la idea principal
%          de cada uno, por decision del autor).
%
% ⚠️ Reglamento art. 11.4.a: «no es original el trabajo que contenga un exceso de
%    informacion de otras fuentes, AUNQUE ESTEN DEBIDAMENTE CITADAS». Defensa de este
%    capitulo: cada trabajo se trae por lo que decide algo en los capitulos 5 a 9, y se
%    dice explicitamente que decide.
%
% ⚠️ Excluidos por no citables: Wang/Zhengzhou, Stirbu, Al Farib.
% =====================================================================================

Este capítulo recorre lo que ya existe para corregir el efecto del ruido, ordenado por
\textbf{cuándo} actúa cada familia de métodos respecto a la ejecución del circuito. Se ha seleccionado
este criterio porque es el que separa a este trabajo de todo lo demás.

\section{La mitigación clásica y su precio}
\label{sec:qem_clasico}

La corrección de errores completa (QEC) exige cientos de qubits físicos por cada qubit
lógico, inviable para las máquinas actuales \cite{preskill2018nisq}. La
mitigación de errores (QEM) es la alternativa: no elimina el ruido, sino que
estima el resultado que se habría obtenido sin él \cite{cai2023quantum}.

Las técnicas de referencia son la \textbf{extrapolación a ruido cero} (ZNE), que trata de estimar cuál habría sido el resultado de un circuito si el ruido del hardware no estuviera presente \cite{temme2017error}; la \textbf{cancelación probabilística de errores} (PEC), que utiliza métodos estadísticos para compensar los errores introducidos durante la ejecución del circuito, aunque requiere un mayor número de ejecuciones \cite{cai2023quantum}; y la \textbf{mitigación de lectura} (M3), centrada en corregir los errores que aparecen al medir el estado final de los qubits \cite{nation2023suppressing}. 

Sus planteamientos son distintos, pero comparten los dos rasgos
que este trabajo toma como punto de partida: todas \textbf{consumen ejecuciones adicionales}
del procesador, el recurso más escaso, y todas actúan \textbf{después} de ejecutar.

\section{Modelos predictivos}
\label{sec:ml_qem}

Una forma de evitar el elevado coste de las técnicas tradicionales de mitigación consiste en sustituir las ejecuciones extra por un modelo entrenado.
\citeA{liao2024machine} demostraron que una red puede aprender la corrección directamente a
partir del valor ruidoso, y validaron este enfoque en hardware IBM real hasta un centenar de
qubits. De sus resultados, sale un dato que este trabajo aprovecha: entre los modelos tabulares,
\textbf{Random Forest es la mejor referencia}, y por eso entra en la comparativa del
Capítulo~\ref{cap:comparativa}.

Sobre esa idea pero en la línea de grafos, GTraQEM \cite{bao2025gtraqem} representa el circuito como
un DAG (\textit{Directed Acyclic Graph}) y lo procesa con un \textit{graph transformer} apoyado en un \textbf{nodo virtual}
conectado a todos los demás, que recoge el contexto global. Esta idea del nodo virtual es una de las características adoptadas en este trabajo, aunque no se emplea la misma arquitectura por los motivos expuestos en el Capítulo~\ref{cap:comparativa}. QEMFormer \cite{bao2025qemformer} se encuentra entre los métodos más avanzados y combina información general del circuito con detalles sobre su estructura.

\textbf{Aquí se encuentra una diferencia importante en la memoria con respecto a esos trabajos:} todos ellos toman el \textbf{valor
ruidoso medido como entrada}, de modo que hay que haber ejecutado el circuito para poder
corregirlo. Además, los resultados de \citeA{placidi2026deep} indican que esta información es la que más contribuye al rendimiento del modelo y sin ella el modelo reduce considerablemente su precision.

\section{Técnicas pre-ejecución}
\label{sec:preejecucion}

Existe una familia distinta que decide \textbf{antes} de gastar la máquina. La herramienta
oficial de IBM, mapomatic \cite{nation2023suppressing}, analiza distintas formas de asignar un mismo circuito al hardware disponible y selecciona la que previsiblemente producirá menos errores; es una
heurística analítica, sin aprendizaje. \citeA{hartnett2024learning} dan un paso más implementando técnicas de aprendizaje
automático con datos medidos en hardware real.

Un linaje reciente sí predice magnitudes sobre el grafo del circuito antes de ejecutarlo, y es
lo más cercano a esta propuesta. Entre estos trabajos, QuEst \cite{wang2025quest} estima la fiabilidad esperada de un circuito, mientras que \citeA{liu2026output} utilizan una red basada en grafos para predecir directamente el resultado que producirá.  
SQuASH \cite{martyniuk2025benchmarking} compara una red de grafos con Random Forest sobre los mismos circuitos, \textbf{precedente directo de esta comparativa}. Y \citeA{tudisco2025graph}, antes de
proponer la suya, demuestran que el número de qubits, la profundidad y el recuento de puertas no bastan para separar los casos: hay que meter el circuito entero en el modelo en vez
de resumirlo.

Ninguno de los cinco predice la degradación de un observable, y ninguno modela el paso del tiempo.

\section{El hardware cambia con el tiempo}
\label{sec:drift}

\citeA{hirasaki2023detection} midieron que las tasas de error no se degradan de forma progresiva,
sino a saltos bruscos que aparecen en cualquier momento, y que la recalibración diaria
no basta para capturarlos. \citeA{huo2026anchor} lo confirman por otro lado: el mismo
programa ejecutado en momentos distintos devuelve resultados diferentes. Esto 
justifica que el conjunto de datos de esta memoria tenga un eje temporal. 


\section{El hueco en la literatura}
\label{sec:gap}

% ⚠️ FORMATO: la plantilla oficial pone el \caption DEBAJO. Se sigue la plantilla.
% ⚠️ Va anclada con [t] y en \small A PROPOSITO: con [htbp] y cuerpo normal no cabia al
%    pie de su pagina, flotaba a la siguiente y arrastraba el ultimo item de la lista,
%    dejando dos tercios de pagina en blanco y sacando el capitulo de sus tres paginas.
% Las cuatro columnas son los cuatro diferenciadores enunciados debajo de la tabla.
% ⚠️ Si alguna fila cambia, cambiarla tambien en doc/SoTA/comparative_analysis.md.
\begin{table}[t]
    \centering
    \small
    \begin{tabular}{@{}lcccc@{}}
        \toprule
        \textbf{Trabajo} & \textbf{Pre-ejec.} & \textbf{Magnitud} & \textbf{Multi-salida} & \textbf{Deriva} \\
        \midrule
        ZNE, PEC, M3                            & no & sí & no & no \\
        \citeNP{liao2024machine}                & no & sí & no & parcial \\
        GTraQEM, QEMFormer                      & no & sí & sí & no \\
        mapomatic, \citeNP{hartnett2024learning}& sí & no & no & no \\
        \citeNP{wang2025quest}                  & sí & sí & no & no \\
        \citeNP{liu2026output}                  & sí & sí & no & no \\
        \citeNP{aktar2024graph}                 & sí & sí & no & no \\
        \citeNP{martyniuk2025benchmarking}      & sí & sí & no & no \\
        \citeNP{tudisco2025graph}               & sí & no & no & no \\
        \midrule
        \textbf{Este trabajo}                   & \textbf{sí} & \textbf{sí} & \textbf{sí} & \textbf{sí} \\
        \bottomrule
    \end{tabular}

    \caption{El estado del arte frente a los cuatro rasgos que definen esta propuesta.}
    \label{tab:hueco}

    {\footnotesize Fuente: elaboración propia.}
\end{table}

La Tabla~\ref{tab:hueco} resume el corpus frente a los cuatro rasgos que definen este trabajo.

Ninguna fila reúne las cuatro columnas: 

\begin{enumerate}
    \item \textbf{Pre-ejecución} (la entrada del modelo nunca incluye el resultado medido).
    \item  \textbf{Magnitud}, que es lo único que permite corregir y no solo descartar.
    \item  La \textbf{multi-salida} sobre varios observables.
    \item El \textbf{paso del tiempo como eje }, con días de calibración reales y disjuntos entre entrenamiento y prueba.
\end{enumerate}
